\begin{textsample}{POS Dim 1 – ap – Score 43.00 – ap/ap\_ef\_9.txt}  \label{ex:f1_pos_173}
4.3.1.3.
Educação Física > Movimento é, assim, \textbf{uma} ação em \textbf{que} \textbf{um} \textbf{sujeito}, pelo seu"se-movimentar", se introduz no Mundo de forma dinâmica e através desta ação percebe e realiza os sentidos/significados em e para o seu meio. ( Trebels apud Kunz, 1991, p. 163 ).
A Educação Física é o componente curricular \textbf{que} tematiza as práticas corporais em suas diversas formas de codificação e significação social, entendidas como manifestações das possibilidades \textbf{expressivas} dos \textbf{sujeitos}, produzidas por diversos grupos sociais no decorrer da história.
Nessa concepção, o movimento humano está sempre inserido no âmbito da cultura e não se limita a \textbf{um} deslocamento espaço-temporal de \textbf{um} segmento corporal ou de \textbf{um} corpo todo.
Em princípio, todas as práticas corporais podem ser objeto do trabalho pedagógico em qualquer etapa e modalidade de ensino.
Ainda assim, alguns critérios de progressão do conhecimento devem ser atendidos, tais como os elementos específicos das diferentes práticas corporais, as características dos \textbf{sujeitos} e os contextos de atuação, sinalizando tendências de organização dos conhecimentos.
Nesses termos, os conteúdos de ensino podem ser considerados como elementos fundamentais para as \textbf{disciplinas}, pois, é por meio deles \textbf{que} se materializa o processo educativo.
Para Vasconsellos ( 2011 ) : > Os conteúdos têm a \textbf{ver} com a exigência de equipar o \textbf{educando} para produção de sentido, com a problemática do sentido para a própria vida ; daí a busca de conteúdos \textbf{que} ajudem o aluno a se localizar, a se posicionar ( mediação semiótica ), usufruir da cultura e a intervir no mundo ( VASCONSELLOS, 2011, p. 162 ).
Dada a condição legalmente instituída para a Educação Física ao ser considerada como componente curricular obrigatório, a partir da LDBEN no. 9.394/96 e do Decreto Lei no 10793/03, art. 1o ( BRASIL, 2003 ), tendo por finalidade nortear a práxis docente na realização de ações \textbf{que} oportunize ao aluno “ o desenvolvimento de suas potencialidades nos aspectos cognitivo, motor, afetivo e social, objetivando seu aprimoramento como seres \textbf{humanos}, excluindo a seletividade e a hiper-competitividade ” ( Resolução no 022 /2010- CEE/AP, Art.2 ), possibilitando o uso crítico e autônomo das variadas manifestações da cultura corporal de movimento.
Dessa forma, a expectativa desse componente curricular e sua reflexão sobre seu objeto de estudo, colabora para a afirmação das camadas \textbf{populares}, quanto a sua inclinação para o diálogo com saberes produzido pela \textbf{humanidade} e na análise sobre a sociedade e sua liberdade de expressão em todas as linguagens, em relação ao \textbf{homem}, vida e mundo, promovendo a autonomia, “ negando a dominação e submissão do \textbf{homem} pelo \textbf{homem} ” ( Coletivo de Autores, 1992, p. 40 ).
A materialidade corpórea deve ser \textbf{uma} discussão simbólica oriunda da antropologia social e com \textbf{um} olhar na antropologia interpretativa de Clifford Geertz. > Antropologia interpretativa é a leitura das sociedades \textbf{enquanto} textos ou como análoga a textos [... ] todos os elementos da cultura analisada devem por tanto ser entendidos à luz desta textualidade, imanente à realidade cultural. ( CHAVES, 2014, p.1 ) \textbf{Contudo} não se deve fechar apenas na questão \textbf{antropológica}, pois segundo Daólio ( 2004, p.41 ), “ a discussão de cultura estaria libertando na Educação Física os chamados elementos da ordem ”.
Os indivíduos precisam ser respeitados em sua história, subjetividade, para atingir \textbf{uma} transformação permissível com elementos ditos de desordem \textbf{que} se distinguem por não ser algo pronto e acabado, mas sim produzidos coletivamente, significados e ressignificados considerando o contexto mundial, global e local.
Sobretudo, se deve prevalecer o bom senso, pois a anamnese é essencial para \textbf{que} se assegure \textbf{um} \textbf{direcionamento} \textbf{que} fortaleça os saberes \textbf{que} formam o universo do ser aprendente.
Portanto, a perspectiva da Educação Física no contexto escolar é compreendida pelo Coletivo de Autores ( 1992 ) como \textbf{uma} prática pedagógica \textbf{que} tematiza formas de atividades \textbf{expressivas} corporais como : jogo, esporte, dança e ginástica, formas estas \textbf{que} \textbf{configuram} \textbf{uma} área de conhecimento \textbf{que} podemos chamar de cultura corporal.
Na perspectiva de Gaya, 2009, a Educação Física escolar, > [... ] constitui -se na \textbf{disciplina} \textbf{que} no interior da escola trata da Cultura Corporal do Movimento Humano.
Como tal cabe a Educação Física no espaço escolar intervir na criação, configuração e modelação de sentidos para as diversas manifestações da Cultura Corporal do Movimento Humano. ( p.25 ) Acompanhando as referidas tendências, apresenta -se a Educação Física Escolar no Sistema educacional do estado do Amapá como \textbf{uma} \textbf{vertente} pedagógica da cultura corporal de movimento.
A Educação Física nos anos iniciais do Ensino Fundamental possibilita aos alunos desenvolver -se integralmente, preparando -o para o aprendizado sobre a corporeidade, a produção de conhecimento e as ressignificações das práticas corporais de acordo com o texto em \textbf{que} o aluno está inserido, dentro de \textbf{uma} sequência didática e \textbf{metodológica} \textbf{que} proporcionará \textbf{um} aprofundamento da base do conhecimento apreendido nos anos iniciais do Ensino Fundamental.
De forma, \textbf{que} o aluno tenha assegurado o direito de aprendizagem dentro do componente curricular e \textbf{que} este se conecte aos outros componentes formando assim, \textbf{uma} teia de conhecimento indissociável.
Assim, o presente documento busca democratizar, humanizar, contextualizar e diversificar o ensino da Educação Física Escolar, saindo das visões biologicista, tecnicista e esportivista para \textbf{uma} concepção mais abrangente \textbf{que} contemple todas as dimensões humanas envolvidas nas manifestações da cultura corporal do movimento.
O objetivo da organização de este currículo é de proporcionar discussão e reflexão sobre a prática corporal para além da sala de aula, ampliando todo o conhecimento \textbf{que} o aluno tem e a ressignificação das suas práticas na perspectiva de \textbf{instigar}, provocar, sensibilizar o aluno para o desejo de aprender e colaborar de forma protagonista dentro do seu ambiente de aprendizagem.
Desta forma, o currículo de Educação Física está \textbf{baseado} na \textbf{pedagogia} \textbf{histórico-crítica} e na \textbf{psicologia} \textbf{histórico-cultural}.
O trato do conhecimento segundo os PCN‟s ( 1997, p.28 ) apresenta o seguinte enfoque : “ Independente de qual seja o conteúdo escolhido, os processos de ensino e aprendizagem devem considerar as características dos alunos em todas as dimensões ( cognitiva, corporal, afetiva, ética, estética, de relação interpessoal e inserção social ) ”.
Além das considerações elencadas, é necessário também \textbf{que} sejam considerados os conhecimentos \textbf{que} o aluno traz de sua realidade cotidiana, devendo ser contextualizados e complementados por conhecimentos mais elaborados na escola, tornando as aulas de Educação Física mais ricas, variadas, significativas e prazerosas.
Nesta conjuntura, espera -se \textbf{que} o \textbf{educando} seja o protagonista no processo de ensino aprendizagem, construindo suas competências por meio da \textbf{problematização} dos temas da cultura corporal de movimento nas dimensões \textbf{conceitual}, procedimental e atitudinal de maneira espontânea, autônoma e significativa e livre da seletividade, da priorização ao aprimoramento das técnicas dos movimentos esportivizados, do desempenho das capacidades físicas e das habilidades motoras, e por fim, da hiper-competitividade.
Conforme a LDB ( 1996 ), a Educação Básica se objetiva em proporcionar \textbf{uma} formação básica para a cidadania.
Essa concepção também alicerça a Base Nacional Comum Curricular, \textbf{que} define quais são as aprendizagens essenciais \textbf{que} todos os alunos têm o direito de adquirir ao longo da educação básica, ela está orientada pelos princípios éticos, estéticos e políticos \textbf{que} visam a formação humana em suas múltiplas dimensões e a construção de \textbf{uma} sociedade \textbf{justa}, democrática e inclusiva.
A BNCC tem como premissa \textbf{uma} educação integral, \textbf{que} visa o pleno desenvolvimento do estudante, seu crescimento como cidadão e sua qualificação para o trabalho.
Na versão homologada da BNCC, os organizadores expressam as aprendizagens essenciais em dez competências gerais, elas definem o cidadão \textbf{que} queremos formar e norteiam a educação \textbf{que} queremos para todos.
Na proposta da Base Curricular, os componentes curriculares, entre eles a educação física, têm o papel de tematizar as práticas corporais em suas diversas formas de codificação e significação social, abordando -as como fenômeno cultural dinâmico, diversificado, pluridimensional, singular e contraditório, desenvolvendo autonomia para a apropriação e utilização da cultura corporal de movimento, favorecendo sua participação de forma confiante e autoral na sociedade, conforme está apresentado nos objetivos de aprendizagem.
Para isso, os conhecimentos das áreas estão mobilizados não só para entender ou explicar a realidade, como também para fazer escolhas a partir desse entendimento e agir em \textbf{uma} determinada \textbf{direção}, assim reforçando a importância dos processos cognitivos : \textbf{percepção}, atenção, memória, \textbf{raciocínio}, procedimentos em contextos de investigação e criação de soluções.
Para Paim e Bonorino ( 2009, p.? ), “ Atualmente, a área da Educação Física evoluiu de tal forma \textbf{que} abrange múltiplos conhecimentos produzidos e usufruídos pela sociedade em geral [... ] ”.
Esta afirmativa resume os argumentos \textbf{que} justificam a presença deste componente no currículo escolar.
Esta reelaboração busca propor \textbf{uma} ressignificação do currículo, comprometendo -o com a formação integral do \textbf{educando}.
Considerando os pressupostos para o ensino da Educação Física, são apresentadas as competências específicas desse componente para o ensino fundamental, as unidades temáticas, objetos de Conhecimento e habilidades a serem desenvolvidas no decorrer dos Nove Anos do Ensino Fundamental.
É importante \textbf{salientar} \textbf{que} a organização das unidades temáticas se \textbf{baseia} na compreensão de \textbf{que} o caráter lúdico está presente em todas as práticas corporais, ainda \textbf{que} essa não seja a finalidade da Educação Física na escola.
Ao brincar, dançar, jogar, praticar esportes, ginásticas ou atividades de aventura, para além da \textbf{ludicidade}, os estudantes se apropriam das lógicas intrínsecas ( regras, códigos, rituais, sistemáticas de funcionamento, organização, táticas etc. ) a essas manifestações, assim como trocam entre si e com a sociedade as representações e os significados \textbf{que} lhes são atribuídos.
Por essa razão, a delimitação das habilidades privilegia oito dimensões de conhecimento : Dimensões do conhecimento 1.
I – Experimentação : refere -se à dimensão do conhecimento \textbf{que} se origina pela vivência das práticas corporais, pelo envolvimento corporal na realização das mesmas.
São conhecimentos \textbf{que} não podem ser acessados sem passar pela vivência corporal, sem \textbf{que} sejam efetivamente experimentados.
Trata -se de \textbf{uma} possibilidade única de apreender as manifestações culturais tematizadas pela Educação Física e do estudante se perceber como sujeito “ de carne e osso ”.
Faz parte dessa dimensão, além do \textbf{imprescindível} acesso à experiência, cuidar para \textbf{que} as sensações geradas no momento da realização de \textbf{uma} determinada vivência sejam positivas ou, pelo menos, não sejam desagradáveis a ponto de gerar rejeição à prática em si. 2.
II – Uso e apropriação : refere -se ao conhecimento \textbf{que} possibilita ao estudante ter condições de realizar de forma autônoma \textbf{uma} determinada prática corporal.
Trata -se do mesmo tipo de conhecimento gerado pela experimentação ( saber fazer ), mas dele se diferencia por possibilitar ao estudante a competência necessária para potencializar o seu envolvimento com práticas corporais no lazer ou para a saúde. \textbf{Diz} respeito àquele rol de conhecimentos \textbf{que} viabilizam a prática efetiva das manifestações da cultura corporal de movimento não só durante as aulas, como também para além delas. 3.
III - Fruição : \textbf{implica} a apreciação estética das experiências sensíveis geradas pelas vivências corporais, bem como das diferentes práticas corporais oriundas das mais diversas épocas, lugares e grupos.
Essa dimensão está vinculada com a apropriação de \textbf{um} conjunto de conhecimentos \textbf{que} permita ao estudante desfrutar da realização de \textbf{uma} determinada prática corporal e/ou apreciar essa e outras tantas quando realizadas por outros. 4.
IV – Reflexão sobre a ação : refere -se aos conhecimentos originados na observação e na análise das próprias vivências corporais e daquelas realizadas por outros.
Vai além da reflexão espontânea, gerada em toda experiência corporal.
Trata -se de \textbf{um} ato intencional, orientado a formular e empregar estratégias de observação e análise para : ( a ) resolver desafios peculiares à prática realizada ; ( b ) apreender novas modalidades ; e ( c ) adequar as práticas aos interesses e às possibilidades próprios e aos das pessoas com quem compartilha a sua realização. 5.
V - Construção de valores : vincula -se aos conhecimentos originados em discussões e vivências no contexto da tematização das práticas corporais, \textbf{que} possibilitam a aprendizagem de valores e normas voltadas ao exercício da cidadania em prol de \textbf{uma} sociedade democrática.
A produção e \textbf{partilha} de atitudes, normas e valores ( positivos e negativos ) são \textbf{inerentes} a qualquer processo de socialização.
No entanto, essa dimensão está diretamente associada ao ato intencional de ensino e de aprendizagem e, portanto, \textbf{demanda} intervenção pedagógica orientada para tal fim.
Por esse motivo, a BNCC se concentra mais especificamente na construção de valores relativos ao respeito às diferenças e no combate aos preconceitos de qualquer natureza.
Ainda assim, não se pretende propor o tratamento apenas desses valores, ou fazê -lo só em determinadas etapas do componente, mas assegurar a superação de estereótipos e preconceitos expressos nas práticas corporais. 6. \textbf{VI} – Análise : está associada aos conceitos necessários para entender as características e o funcionamento das práticas corporais ( saber sobre ).
Essa dimensão reúne conhecimentos como a classificação dos esportes, os sistemas táticos de \textbf{uma} modalidade, o efeito de determinado exercício físico no desenvolvimento de \textbf{uma} capacidade física, entre outros. 7.
VII - Compreensão : está também associada ao conhecimento \textbf{conceitual}, mas, diferentemente da dimensão anterior, refere -se ao esclarecimento do processo de inserção das práticas corporais no contexto sociocultural, reunindo saberes \textbf{que} possibilitam compreender o lugar das práticas corporais no mundo.
Em linhas gerais, essa dimensão está relacionada a temas \textbf{que} permitem aos estudantes interpretar as manifestações da cultura corporal de movimento em relação às dimensões éticas e estéticas, à época e à sociedade \textbf{que} as gerou e as modificou, às razões da sua produção e transformação e à vinculação local, nacional e global.
Por exemplo, pelo estudo das condições \textbf{que} permitem o surgimento de \textbf{uma} determinada prática corporal em \textbf{uma} dada região e época ou os motivos pelos quais os esportes praticados por \textbf{homens} têm \textbf{uma} visibilidade e \textbf{um} tratamento midiático diferente dos esportes praticados por mulheres. 8.
VIII - Protagonismo comunitário : refere -se às atitudes/ações e conhecimentos necessários para os estudantes participarem de forma confiante e autoral em decisões e ações orientadas a democratizar o acesso das pessoas às práticas corporais, tomando como referência valores favoráveis à convivência social.
Contempla a reflexão sobre as possibilidades \textbf{que} eles e a comunidade têm ( ou não ) de acessar \textbf{uma} determinada prática no lugar em \textbf{que} moram, os recursos disponíveis ( públicos e privados ) para tal, os agentes envolvidos nessa configuração, entre outros, bem como as iniciativas \textbf{que} se dirigem para ambientes além da sala de aula, orientadas a interferir no contexto em busca da materialização dos direitos sociais vinculados a esse universo.
Vale ressaltar \textbf{que} não há nenhuma \textbf{hierarquia} entre essas dimensões, tampouco \textbf{uma} ordem necessária para o desenvolvimento do trabalho no âmbito didático.
Cada \textbf{uma} delas \textbf{exige} diferentes \textbf{abordagens} e graus de complexidade para \textbf{que} se tornem relevantes e significativas.
Considerando as características dos conhecimentos e das experiências próprias da Educação Física, é importante \textbf{que} cada dimensão seja sempre abordada de modo integrado com as outras, levando -se em conta sua natureza vivencial, experiencial e subjetiva.
Assim, não é possível \textbf{operar} como se as dimensões pudessem ser tratadas de forma isolada ou sobreposta.
Considerando esses pressupostos, e em articulação com as competências gerais da BNCC e as competências específicas da área de Linguagens, o componente curricular de Educação Física deve garantir aos alunos o desenvolvimento de competências específicas.
No quadro a seguir, pode -se verificar as 10 ( dez ) competências \textbf{que} \textbf{permeiam} o Componente Curricular Educação Física, é possível perceber, como \textbf{dito} anteriormente, nenhuma competência se sobrepõe a outra, se faz necessário \textbf{que} ao final do Ensino Fundamental as 10 competências específicas do Componente Educação Física sejam contempladas ao longo dos 9 anos desta etapa de ensino.

% matched lemmas: abordagem, antropológico, basear, conceitual, configurar, contudo, demandar, direcionamento, direção, disciplina, dizer, educar, enquanto, exigir, expressivo, hierarquia, histórico-crítico, histórico-cultural, homem, humanidade, humanos, implicar, imprescindível, inerente, instigar, justo, ludicidade, metodológico, operar, partilha, pedagogia, percepção, permear, popular, problematização, psicologia, que, raciocínio, salientar, sujeito, um, ver, vertente
\end{textsample}
